\subsection{Comparison with the continuous profile: the system and its residual}\label{sec:P05}
\paragraph{Initial data.}
Sections~\ref{sec:P01} and~\ref{sec:P02} give, for the canonical maximizer $\pi$ and
$n\ge n_0$, where the threshold is that of Section~\ref{sec:P02},
\[
 c\le\rho^*_{ij}:=n\pi_{ij}\le M,\qquad h_D^*,h_A^*\ge\eta>0.
                                                        \tag{TR1}
\]
For the continuous profile of mass one, we use only its factorization, its equations,
and its positivity and Lipschitz bounds (Section~\ref{sec:preliminares}):
\[
 \nu_1(r,c)=b_0 r_R(r)r_C(c)u_D(r-c)u_A(r+c),\quad
 c_\nu\le\nu_1\le C_\nu,\quad \kappa_\nu\le u_D,u_A\le1.
                                                        \tag{TR2}
\]
Its row and column marginals equal one, and $u_f+T_f\nu_1=1$ for $f=D,A$, where $T_f\nu_1=m_f$ is
the projection of $\nu_1$ onto family $f$.
Recall that $L_\nu$ is a Lipschitz constant for $\nu_1$ with respect to the maximum metric.

\paragraph{Grid, operators, and gauge.}
Set $x_i=(i-1/2)/n$. The lines $D_k$ and $A_l$ have parameters
$k/n$ and $l/n$, with $k=i-j$ and $l=i+j-1$ at square $(i,j)$.
The four families consist of $n,n,2n-1,2n-1$ lines, respectively. Define
\[
 (\mathcal R_f p)_\ell=\frac1n\sum_{(i,j)\in\ell}p_{ij},\qquad
 (\mathcal T z)_{ij}=z_{R,i}+z_{C,j}+z_{D,i-j}+z_{A,i+j-1}.
                                                        \tag{TR3}
\]
Write $\mathcal R=(\mathcal R_R,\mathcal R_C,
\mathcal R_D,\mathcal R_A)$. The inner products are
$\langle p,q\rangle_\square=n^{-2}\sum pq$ and
$\langle h,k\rangle_\ell=n^{-1}\sum_{f,\ell}h_{f,\ell}k_{f,\ell}$,
so $\mathcal R^*=\mathcal T$. We use the same space for the domain and range:
\[
 E_n=\{z:\ \sum_i z_{R,i}=\sum_j z_{C,j}\}=v_0^\perp,\qquad
 v_0=(\mathbf1_R,-\mathbf1_C,0_D,0_A).
                                                        \tag{TR4}
\]
Every potential can be moved into $E_n$ by subtracting $\tau v_0$, where
$\tau=(\overline z_R-\overline z_C)/2$. This leaves $\mathcal Tz$ unchanged.

\paragraph{The small root and the exact system.}
Let $\psi_n(t)=t(1-t/n)^3$. Since
$\psi_n'(t)=(1-t/n)^2(1-4t/n)>0$ for $0<t<n/4$,
its inverse on this branch, denoted by $\chi_n$, is unique and smooth.
When the arguments lie in the image of this branch, set
\[
 \rho(z)_{ij}=\chi_n(e^{(\mathcal Tz)_{ij}}),\qquad
 u_D(z)=e^{z_D},\quad u_A(z)=e^{z_A},
\]
\[
 F_n(z)=(\mathcal R_R\rho(z)-1,\mathcal R_C\rho(z)-1,
 u_D(z)+\mathcal R_D\rho(z)-1,u_A(z)+\mathcal R_A\rho(z)-1).
                                                        \tag{TR5}
\]
We always have $F_n(z)\in E_n$. The sums of its R and C components
are exactly equal because both sum to $n^{-1}\sum_{ij}\rho(z)_{ij}-n$.

\paragraph{The canonical maximizer as a solution of the system.}
By Section~\ref{sec:P02}, no cap is active and all the multipliers $\mu_{ij}$ in (P1.3) vanish.
Before fixing the gauge, take
\[
 z^*_{R,i}=\tfrac12\log n-\alpha_i-1,\quad
 z^*_{C,j}=\tfrac12\log n-\beta_j-1,\quad
 z_D^*=\log h_D^*,\quad z_A^*=\log h_A^*.
\]
With this choice, (P1.3) gives $\psi_n(\rho^*_{ij})=e^{(\mathcal Tz^*)_{ij}}$.
For large $n$, we have $\rho^*/n<1/4$ by (TR1). The uniqueness of the small
root gives $\rho(z^*)=\rho^*$ and $F_n(z^*)=0$.
The gauge adjustment preserves these identities.
Thus this solution of the system represents the canonical maximizer itself. It does not
define a distinct branch of solutions.

\paragraph{Sampling the continuous profile and quadrature on all fibers.}
Before fixing the same gauge, set
\[
 \begin{aligned}
 z^0_{R,i}&=\tfrac12\log b_0+\log r_R(x_i), & z^0_{C,j}&=\tfrac12\log b_0+\log r_C(x_j),\\
 z^0_{D,k}&=\log u_D(k/n), & z^0_{A,l}&=\log u_A(l/n).
 \end{aligned}                                          \tag{TR6}
\]
Then $e^{(\mathcal Tz^0)_{ij}}=\widehat\nu_{ij}:=\nu_1(x_i,x_j)$.
Let $Q=\max\{2C_\nu,M\}$ and fix $n\ge\max\{n_0,\lceil8Q\rceil\}$.
Since $\psi_n(2C_\nu)\ge2C_\nu(7/8)^3>C_\nu$,
the initial root exists and satisfies
\[
 c_\nu\le\rho(z^0)\le2C_\nu,\qquad
 0\le\rho(z^0)-\widehat\nu\le12C_\nu^2/n.
                                                        \tag{TR7}
\]
The last inequality uses $t-\psi_n(t)\le3t^2/n$.
In the rows and columns, the squares are represented by the midpoints of intervals of length $1/n$.
On $D_k$, the row coordinate ranges over
$[\max(0,k/n),\min(1,1+k/n)]$, and on $A_l$ it ranges over
$[\max(0,l/n-1),\min(1,l/n)]$.
The endpoints are multiples of $1/n$, and the values of $x_i$ encountered are
exactly the midpoints of the corresponding partition.
Each restriction of $\nu_1$ has Lipschitz constant at most $L_\nu$.
Integration is with respect to the row or column coordinate, with no factor of $\sqrt2$.
The error per subinterval is at most $L_\nu/(4n^2)$, and there are at most
$n$ subintervals. Uniformly, even on the shortest diagonals,
\[
 |\mathcal R_f\widehat\nu-T_f\nu_1(\text{parameter})|
 \le L_\nu/(4n),\qquad
 \|F_n(z^0)\|_\infty\le C_0/n,\quad C_0=12C_\nu^2+L_\nu/4.
                                                        \tag{TR8}
\]

\paragraph{Control along the segment without assuming closeness.}
On the segment $z(t)=(1-t)z^0+tz^*$, the exponential $e^{\mathcal Tz(t)}$ interpolates geometrically
between the arguments of the two roots. By the monotonicity of $\chi_n$, each
$\rho(z(t))_{ij}$ lies between its two endpoint values. Similarly, the
exponential vacancies interpolate geometrically between their endpoint values. Thus
\[
 \begin{gathered}
 a\le\rho(z(t))\le Q,\quad \rho(z(t))/n\le1/8,\qquad
 \zeta\le u_D(z(t)),u_A(z(t))\le1,\\
 a=\min(c_\nu,c),\qquad\zeta=\min(\kappa_\nu,\eta)>0.
 \end{gathered}                                         \tag{TR9}
\]
The intermediate points need not satisfy the marginal constraints.
